\documentclass[10pt]{article}

\usepackage[utf8]{inputenc}

\usepackage[T1]{fontenc}

\usepackage{amsmath}

\usepackage{amsfonts}

\usepackage{amssymb}

\usepackage[version=4]{mhchem}

\usepackage{stmaryrd}

\usepackage{bbold}



\begin{document}

\begin{enumerate}

  \item Проблема существования граничных 3 н а ч е н й, то есть вопрос о том, при каких условиях и в каком смысле сушествуют граничные значения $f(z)$ при приближении точки $z$ к Г. Эта проблема, как и последующие, иначе может быть сформулирована как задача о выделении достаточно обширных классов аналитич. функций в $D$, имеющих в том или ином смысле граничные значения на достаточно массивных множествах точек Г.



  \item Проблема граничного представления $f(z)$, то есть вопрос о том, при каких условиях и при помощи какого аналитич. аппарата может быть выражена зависимость функции $f(z)$ от ее граничных значений на $\Gamma$. Здесь, очевидно, для различных классов аналитич. функций аналитич. аппарат будет варьироваться.



  \item П робле ма е д и н с т в е н ности, или вопрос о том, какими свойствами должно обладать множество $E \subset \Gamma$, чтобы две аналитич. функции того или иного класса совпадали всюду в $D$, если их граничные значения на $E$ совпадают.



\end{enumerate}



Лum.: [1] Математическая энциклопедия, т. 1, М., 1977, кол. 1098-1102; [2] Гол у б е в В. В., Однозначные аналитические функции. Автоморфные функции, М., 1961; [3] Го л уз и н Г. М., Геометрическая теория функций комплексного переменного, 2 изд., М., 1966; [4] П р в алов И. И., Граничные свойства аналитических функций, 2 изд., М., $1950 . \quad$ Е.Д. Соломенцев. TPACCMAHA AITEБPA, в н еш н я я ал г е б ра,- алгебра, порожденная антикоммутирующими образующими $x_{1}, x_{2}, \ldots, x_{n}$, то есть совокупность всевозможных линейных комбинаций из произведений образующих $x_{i}$, в которых все сомножители различны, так как



\[

x_{i} x_{k}+x_{k} x_{i}=0

\]



и, в частности, $x_{k}^{2}=0$ при любом $k$. Размерность Г.а. как линейного пространства равна $2^{n}$, базис состоит из $2^{n}$ одночленов:



\[

\begin{gathered}

1, x_{i}, i \leqslant n, x_{i} x_{j}, i<j \leqslant n, \\

x_{i} x_{j} x_{k}, i<j<k \leqslant n, \ldots, x_{1} x_{2} \ldots x_{n} .

\end{gathered}

\]



Любой элемент Г. а. $f(x)$ можно представить в виде следующей конечной суммы:



\[

f(x)=f+\sum_{i} f^{i} x_{i}+\sum_{i<j} f^{i j} x_{i} x_{j}+\ldots+f^{1 \ldots n} x_{1} x_{2} \ldots x_{n} .

\]



На случай грассмановых переменных обобшается ряд понятий обычного анализа, в частности дифференцирование и интегрирование. Чтобы найти левую производную от одночлена $x_{1} \ldots x_{\beta}$ по переменной $x_{\alpha}$, нужно, пользуясь (1), переставить $x_{\alpha}$ на первое слева место и вычеркнуть ее. Аналогично определяется правая производная. Производная от общего элемента Г. а. есть сумма производных от одночленов в разложении (2). Интеграл на Г.а. задается правилами Березина:



\[

\int d x^{\alpha}=0, \quad \int d x^{\alpha} x_{\beta}=\delta_{\beta}^{\alpha}

\]



при этом кратный интеграл понимается как повторный. Символ $d x^{\alpha}$ не есть обычный дифференциал, его следует трактовать формально. Интеграл на Г. а. обладает нек-рыми свойствами обычного интеграла, в частности возможно интегрирование по частям. С другой стороны, интегрирование на Г. а. эквивалентно дифференцированию:



\[

\int d x^{\alpha} f(x)=\partial f(x) / \partial x_{\alpha}

\]



Интегрирование по грассмановым переменным позволяет построить функциональный интеграл, представляющий Грина функции фермионных полей. ГРАССМАНА ЛАГРАНХЕВО МНОГООБРАЗИЕ $\Lambda_{n}-$ множество всех Лагранжа плоскостей линейного $2 n$-мерного симплектического пространства.



$\Lambda_{1}=\mathbb{R} P^{1} \cdot \Lambda_{2}$ изоморфно квадрике $x^{2}+y^{2}+z^{2}-u^{2}-w^{2}=$ $=0$ в $\mathbb{R} P^{4} ; \Lambda_{2}$ является факторизацией $S^{2} \times S^{1}$ по антипо- дальной инволюции $(x, y) \mapsto(-x,-y) ; \operatorname{dim} \Lambda_{n}=n(n+1) / 2$. Г. л. м. является однородным пространством: $\Lambda_{n}=U_{n} / O_{n}$.



ГPАССМАНОВЫ ПЕРЕМЕННЫЕ - образуюшие Грассмана алгебры $x_{i}$, удовлетворяющие соотношениям



\[

\left\{x_{i}, x_{j}\right\}=x_{i} x_{j}+x_{j} x_{i}=0 .

\]



Произвольный элемент алгебры с $n$ образующими есть многочлен вида



\[

f(x)=\sum_{a_{i}=0,1} C\left(a_{1}, \ldots, a_{n}\right) x_{n}^{a_{n}} \ldots x_{1}^{a_{1}} .

\]



В силу перестановочных соотношений (коммутативности) при $i=j$ получается $x_{i}^{2}=0$, так что степени образующих выше первой исчезают. Произведение $k$ различных Г. П. называется о д н очл е ном с те пе н и $k$. Элементы алгебры Грассмана, представимые в виде линейной комбинации одночленов четной степени, называются ч е т ным и, а представимые в виде линейной комбинации одночленов нечетной степени - н еч е т ны м и. На грассмановой алгебре можно ввести дифференциальное исчисление (см. [1], [2]). Определяются левая $\frac{\partial}{\partial x_{p}} f$ и правая $f \frac{\partial}{\partial x_{p}}$ производные элемента $f$ по Г. П. $x_{p}$. В силу линейности операции дифференцирования для этого достаточно задать производные на одночленах. Пусть



\[

\begin{aligned}

& \frac{\partial}{\partial x_{p}} x_{i_{1}} \ldots x_{i_{s}}=\delta_{i_{1} p} x_{i_{2}} \ldots x_{i_{s}}-\delta_{i_{2} p} x_{i_{1}} x_{i_{3}} \ldots x_{i_{s}}+\ldots \\

& \ldots+(-1)^{s-1} \delta_{i_{s}} p_{i} \ldots x_{i_{s-1}}, \\

& x_{i_{1}} \ldots x_{i_{s}} \frac{\partial}{\partial x_{p}}=\delta_{i_{s}} x_{i_{1}} \ldots x_{i_{s-1}}-\delta_{i_{s-1}} p^{x_{1}} \ldots x_{i_{s-2}} x_{i_{s}}+\ldots \\

& \ldots+(-1)^{s-1} \delta_{i_{1} p} x_{i_{2}} \ldots x_{i_{s}} .

\end{aligned}

\]



Производные обладают следующими свойствами:



\begin{enumerate}

  \item Пусть $f_{1}-$ четный, $f_{2}-$ нечетный, а $f-$ произвольный элементы грассмановой алгебры. Тогда

\end{enumerate}



\[

\begin{aligned}

& \frac{\partial}{\partial x_{p}}\left(f_{1} f\right)=\left(\frac{\partial}{\partial x_{p}} f_{1}\right) f-f_{1}\left(\frac{\partial}{\partial x_{p}} f\right) \\

& \frac{\partial}{\partial x_{p}}\left(f_{2} f\right)=\left(\frac{\partial}{\partial x_{p}} f_{2}\right) f+f_{2}\left(\frac{\partial}{\partial x_{p}} f\right)

\end{aligned}

\]



и аналогичные формулы для левых производных.



\begin{enumerate}

  \setcounter{enumi}{1}

  \item $\frac{\partial}{\partial x_{1}}\left(\frac{\partial}{\partial x_{2}} f\right)=-\frac{\partial}{\partial x_{2}}\left(\frac{\partial}{\partial x_{1}} f\right),\left(\frac{\partial}{\partial x_{1}} f\right) \frac{\partial}{\partial x_{2}}=\frac{\partial}{\partial x_{1}}\left(f \frac{\partial}{\partial x_{2}}\right)$.

\end{enumerate}



Для грассмановой алгебры с четным числом образующих, заданных в виде $x_{1}, x_{1}^{+}, \ldots, x_{n}, x_{n}^{+}$можно ввести операцию и н в ол ю и и, действующую на произвольный элемент вида



\[

\begin{gathered}

f\left(x, x^{+}\right)=\sum_{a_{i}, b_{i}=0,1} C\left(a_{1}, \ldots, a_{n} ; b_{1}, \ldots, b_{n}\right)\left(x_{n}\right)^{a_{n}} \ldots \\

\ldots\left(x_{1}\right)^{a_{1}}\left(x_{n}^{+}\right)^{b_{n}} \ldots\left(x_{1}^{+}\right)^{b_{1}}

\end{gathered}

\]



по правилу



\[

\begin{aligned}

f \rightarrow f^{+}\left(x, x^{+}\right)= & \sum_{a_{i}, b_{i}=0,1} C^{*}\left(a_{1}, \ldots, a_{n} ; b_{1}, \ldots, b_{n}\right)\left(x_{n}\right)^{b_{n}} \ldots \\

& \ldots\left(x_{1}\right)^{b_{1}}\left(x_{n}^{+}\right)^{a_{n}} \ldots\left(x_{1}^{+}\right)^{a_{1}} .

\end{aligned}

\]



Для грассмановых алгебр с инволюцией со скалярным произведением инволютивные Г. П. можно представить в виде $\varepsilon(x), \varepsilon^{+}(x)$, где $x$ пробегает нек-рое множество $M$ с мерой $d x$ (см. [1]). При этом любой элемент грассмановой алгебры представим в виде функционала от функщий с антикоммутируюшими значениями



\[

\begin{gathered}

f\left(\varepsilon, \varepsilon^{+}\right)=\sum_{n, m} \int \varphi\left(x_{1}, \ldots, x_{n} ; y_{1}, \ldots, y_{m}\right) \varepsilon\left(x_{1}\right) \ldots \varepsilon\left(x_{n}\right) \times \\

\times \varepsilon^{+}\left(y_{1}\right) \ldots \varepsilon^{+}\left(y_{m}\right) d x^{n} d y^{m},

\end{gathered}

\]





\end{document}